% !TEX root = ../main.tex
\chapter{Sistemi logici}
\label{chap:logical_systems}

In questo capitolo viene introdotto il formalismo logico necessario per completare il parallelo tra informatica teorica e logica matematica. Dopo una definizione strutturale 
di sistema logico e dei principali paradigmi di deduzione, l'analisi si concentrerà sulla logica intuizionista. Quest'ultima, attraverso l'interpretazione BHK e il concetto 
di inferenzialismo, si distacca dal modello classico per porre la "prova" come elemento centrale della verità. Esploreremo la gerarchia del cubo logico intuizionista, dalla 
logica proposizionale a quella di ordine superiore, definendo le basi sintattiche e semantiche che permetteranno, nel capitolo successivo, di formalizzare l'identità tra 
programmi e dimostrazioni.

\section{Definizione di sistema logico} 
Un sistema logico è formalmente definito dalla terna $\mathcal{L} = \langle \mathcal{S}, \mathcal{V}, \models \rangle$, che storicamente possiede tale significato:
\begin{itemize}[nosep, leftmargin=1.5em]
    \item $\mathcal{S}$ rappresenta la sintassi, ovvero l'insieme delle formule ben formate.
    \item $\mathcal{V}$ indica la semantica (o l'insieme dei valori di verità/modelli).
    \item $\models$ definisce la relazione di conseguenza logica (o soddisfacibilità) tra gli elementi del sistema.
\end{itemize}

% TODO: aggiungere valida, soddisfacibile, contraddizione

\section{Sistemi di deduzione}
Un sistema di sistema di deduzione (o sistema di prova) si definisce come un algoritmo d'inferenza preposto alla generazione di dimostrazioni. Mentre il sistema logico stabilisce, 
sul piano semantico, quali enunciati siano conseguenza logica di altri ($\Gamma \models A$), un sistema di prova definisce un insieme finito di regole sintattiche atte 
a derivare tali conclusioni attraverso passaggi meccanici. In tal senso, il sistema di prova rappresenta il corrispettivo sintattico della relazione di conseguenza logica
$\models$. Tale relazione di derivabilità sintattica viene formalmente indicata con il simbolo di turnstile $\vdash$.

\subsection{Correttezza (Soundness)} 
La correttezza o soundness è la proprita di un sistema di deduzione di derivare solo verita rispetto al sistema logico di riferimento. Nello specifico,
se si prova $A$ nel sistema ($\Gamma \vdash A$), allora $A$ è una verità logica ($\Gamma \models A$). Il mio sistema non produce errori.

\subsection{Completezza (Completeness)}
La completezza o completeness è la proprita di un sistema di deduzione di poter derivare tutte verita possibili rispetto a un sistema logico. Nello specifico,
Se $A$ è una verità logica ($\Gamma \models A$), allora esiste certamente una prova nel mio sistema ($\Gamma \vdash A$). Il sistema di deduzione è abbastanza 
potente da trovare tutte le verità.    

\subsection{Classificazione di sistemi di deduzione}
La classificazione dei sistemi di deduzione è secondo molti punti di vista connfusionaria, tale confuzione è dovuta alla natura algoritmica degli stessi applicata alle filosofie specifiche
di ogni logica. Nonostante i sistemi di deduzione siano indipendenti dalla logiche ne devono rispettare intrinseca definizione di verita derivato dalla semantica. Inoltre sono sistemi d'inferenza
percio algoritmi sintattici, questo loro aspetto viene potenziato e valorizzato in molti sistemi logici mentre in altri viene solamente utilizzato come strumento posteriore. Di seguito sono riportati
i sistemi di deduzione piu utilizzati. 

\subsubsection{Sistemi assiomatici}
I sistemi assiomatici (Sistemi alla Hilbert) sono caratterizzati da un vasto insieme di assiomi logici, verita a priori, e un numero ridotto di regole d'inferenza, solitamente solo \emph{modus ponens}.
In questi sistemi, una dimostrazione è una sequenza lineare fi formule dove ogni passaggio è un assioma o deriva dei precedenti. Sebbene siano fondamentali per lo studio della matematica e delle completezza,
risultano poco intuitivi per la pratica umana poiché richiedono costruzioni logiche molto complesse anche per enunciati semplici.
\begin{center}
    $\begin{array}{rll}
        1. & A \to B & \text{(Assioma o Ipotesi)}       \\
        2. & A        & \text{(Assioma o Ipotesi)}      \\
        3. & B        & \text{(Modus Ponens da 1, 2)}
    \end{array}$
\end{center}

\subsubsection{Deduzione naturale}
Sviluppata da Gerhard Gentzen, la deduzione naturale mira a riprodurre fedelmente il modo in cui il ragionamento matematico avviene "in natura". Il sistema si basa su una struttura simmetrica: per ogni connettivo 
logico esiste una regola di \emph{introduzione} (che definisce come costruire una formula con quel connettivo) e una regola di \emph{eliminazione} (che definisce come utilizzarla). Questo sistema è il pilastro della 
logica intuizionista poiché.
\begin{center}
    \begin{prooftree}
        \hypo{ A \to B }
        \hypo{ A }
        \infer2[($\to E$)]{ B }
    \end{prooftree}
\end{center}

\subsubsection{Calcolo dei sequenti}
Sempre introdotto da Gentzen, il calcolo dei sequenti non opera su singole formule ma su \emph{sequenti} del tipo $\Gamma \vdash \Delta$, che rappresentano una relazione tra un insieme di assunzioni e un insieme di conclusioni. 
A differenza della deduzione naturale, questo sistema eccelle nell'analisi strutturale della prova e nella ricerca automatica delle dimostrazioni. In ambito intuizionista, la distinzione cruciale risiede nella restrizione del lato 
destro del sequente ($\Delta$) a una singola formula, impedendo formalmente la derivazione di teoremi non costruttivi come il terzo escluso
\begin{center}
    \begin{prooftree}
        \hypo{\Gamma \vdash A}
        \hypo{\Gamma, B \vdash C}
        \infer2[($L\to$)]{\Gamma, A \to B \vdash C}
    \end{prooftree}
\end{center}

\subsubsection{Metodo dei tableaux}
Il metodo dei tableaux (o alberi di refutazione) è un sistema di prova basato sulla ricerca della contraddizione. Per dimostrare che $\Gamma \models A$, il sistema tenta di confutare la negazione della conclusione, 
ovvero cerca di mostrare che $\Gamma \cup \{\neg A\}$ è inconsistente. Il sistema procede scomponendo le formule in rami (scelte logiche); se tutti i rami si chiudono in una contraddizione, la formula originale è provata. 
\begin{center}
    $\begin{array}{c}
        A \to B                     \\
        A                           \\
        \neg B                      \\
        \hline
        \swarrow \quad \searrow     \\
        \neg A \quad\quad B         \\
        \otimes \quad\quad \otimes
    \end{array}$
\end{center}

\section{Classificazione dei sistemi logici}
La classificazione dei sistemi logici presenta intrinseche complessita metodologiche. Tali complessita derivano dalla loro natura fondazionale del pensiero matemantico. In un certo senso,
i sistemi logici sono modelli formali del ragionamento, di conseguenza sono modelli di modelli. Percio si puo affermare che sistemi logici definiscono il calcolo del linguaggio. La definizione
di tale calcolo avviene anche essa tramite il linguaggio, questa ricorsivita spesso porta a criticita. Per tale motivazione, il processo di classificazione dei sistemi logici è diviso in piu 
momenti. Le prime formulazioni separavano sintassi e semantica con la conseguenza implicita di scindere la relazione di conseguenza logica dal processo d'inferenza. Le prospettive più moderne, 
tra cui quella categoriale, ricercano una sovrapposizione e un'unificazione di tali concetti.

\subsection{Paradgmi semantici}
Per orientarsi nell'ampio panorama dei sistemi logici, è utile l'individuazione di criteri di confronto rigorosi che fungano da coordinate all'interno dello spazio dei paradigmi semantici.
I criteri utilizzati sono: \emph{la denotazione}, \emph{la natura della relazione di conseguenza} e \emph{l'identità delle dimostrazioni}. Tali criteri permettono di mappare l'evoluzione del 
concetto di verità dal modello classico a quello categoriale. Un'analisi esaustiva di tali criteri richiederebbe tuttavia un apparato formale di cui non disponiamo ancora. 

\subsection{Logica classica}
Il sistema logico classico è emerso storicamente per rispondere alle esigenze di formalizzazione della matematica. Tuttavia, per sua natura la formalizzazione non è stata lineare, portando a una 
struttura che predilige la verità dei risultati rispetto alla dinamica delle dimostrazioni. Una caratteristica cruciale della logica classica è infatti l'assenza di un'interpretazione computazionale 
diretta: in questo contesto, la prova è un mezzo trasparente per stabilire un valore di verità, non un oggetto di calcolo. Per tale ragione, appare metodologicamente più fecondo adottare come base 
la logica intuizionista. Caratterizzata dall'assenza del principio del terzo escluso ($A \vee \neg A$), essa permette una formalizzazione più profonda del concetto di prova, culminando in sistemi 
come la teoria dei tipi semplici. Questi ultimi non solo garantiscono la coerenza del sistema eliminando paradossi logici (come quello di Russell), ma preservano la natura algoritmica della deduzione. 
In questa prospettiva, il sistema logico classico non viene escluso, ma recuperato come estensione: esso si ottiene aggiungendo ai sistemi intuizionisti l'assioma del terzo escluso o la legge della 
doppia negazione ($\neg \neg A \to A$), con la consapevolezza che tale aggiunta rompe la natura strutturale della prova.

% TODO: aggiungere delle richiami alle altre logiche per poi trattare approfonditamente solo la logica intuizionista

\section{Logica intuizionista}

\subsection{Inferenzialismo}
La logica intuizionista rappresenta l'espressione piu coerente dell'inferenzialismo semantico, attraverso l'interpretazione BHK (Brouwer-Heyting-Kolmogorov). Nel inferenzialismo
il significato dei costrutti delle formule non è determinato da condizioni di verita trascendenti, ma è interamente racchiuso nelle loro condizioni di asseribilità, ovvero
nelle regole di introduzione e eliminazione della deduzione naturale. La prospettiva inferenzialista trasforma la logica, in una teoria della giustificazione (\emph{proof-theoretic}) 
dove la verita è la prova stessa anziche una descrizione di stati di verita esterni al sistema logico. La conseguenze di maggior impatto è la fusione  dei concetti di conseguenza logica
e sistema di deduzione, \emph{antirealismo}. 

\subsection{Metaregole intuizioniste}
Una volta definite le basi filosofiche del rapporto con la verità su cui si baseranno le semantiche intuizioniste, 
risulta necessario introdurre le metaregole che permettono di non costruire sistemi inconsistenti:
\begin{itemize}
    \item 
        \textbf{L'armonia}: è il principio fondamentale tra regole di introduzione e eliminazione. 
        Tale principio nega di costruire sistemi in cui l'informazione derivata dalla regola 
        di eliminazione è maggiore dell'informazione inserita dalla regola di introduzione, bilanciamento.
    \item 
        \textbf{La conservatività}: si pone come seconda metaregola cruciale e assicura che l'introduzione di nuovi 
        connettivi non alteri le relazioni logiche già stabilite tra i simboli preesistenti. Un sistema che 
        violi questo principio rischierebbe la trivializzazione, annullando la distinzione tra vero e falso.
\end{itemize}
Per l'intuizionista, queste metaregole non sono meri vincoli sintattici, ma condizioni di possibilità della semantica stessa: esse garantiscono che ogni estensione del linguaggio 
sia "sicura" e che la verità rimanga sempre ancorata a una prova effettiva. Tuttavia, è importante notare che non tutti i sistemi logici aderiscono alla conservatività come requisito 
necessario. Alcuni autori sostengono che l'estensione del linguaggio posso legittimamente espandere il potere deduttivo del sistema, se non viola l'armonia. Per esempio, l'aggiunta
della doppia degazione alla logica intuizionista per ottenere la logica classica rompe la conservatività.

\subsection{Cubo delle logiche intuizioniste}
La classificazione interna alla logica intuizionista puo essere rappresentata attraverso una struttura a cubo definita da tre assi di espansione. Questa classificazione si basa sulla natura delle
entita quantificabili e sulla gerarchia linguaggio. Il sistema logico di partenza, ovvero il vertice meno espressivo, è chiamato \emph{Logica Proposizionale}.
\begin{itemize}
    \item \textbf{Asse $x$ (Individui):}  
        Rappresenta l'espansione verso la \emph{logica del primo ordine intuizionista} (IFOL). Questo asse introduce la capacità di analizzare la struttura interna 
        delle proposizioni attraverso l'uso di variabili individuali e costanti. In questa dimensione, i predicati sono trattati esclusivamente come 
        parametri fissi (costanti sintattiche) e la quantificazione è limitata agli oggetti del dominio.
    \item \textbf{Asse $y$ (Ordini):} 
        Definisce la capacità di quantificare su entità di rango superiore, come proposizioni o proprietà. Questo asse conduce alla \emph{logica di secondo ordine intuizionista}. 
        A differenza dell'asse $x$, qui le proprietà stesse diventano variabili del linguaggio, permettendo di esprimere principi universali sulle classi di oggetti 
        e sulle verità logiche.
    \item \textbf{Asse $z$ (Tipi):}
        Introduce la \emph{gerarchia dei tipi} e l'astrazione funzionale. Questo asse permette di superare la distinzione rigida tra individuo e predicato, 
        organizzando il linguaggio in una struttura ricorsiva di tipi semplici. In questa dimensione, ogni termine possiede un rango definito che ne determina 
        il ruolo di funzione o di argomento.
\end{itemize}
L'intersezione di questi tre assi genera otto sistemi logici distinti, il cui vertice di massima potenza espressiva è rappresentato dalla \emph{logica di ordine superiore intuizionista} (IHOL), 
in cui la distinzione tra formule e termini viene unificata all'interno di una gerarchia tipata completa.

\paragraph{Metodologia semantica.}
Inoltre è importante notare, che almeno per ora, limiteremo la trattazione della semantica delle logiche intuistioniste utilizzando la semantica BHK. Tale semantica sebbene informale (pre-matematica) 
è sia abbastanza precisa che intuitiva da permetterci di cogliere il contenuto informativo della logica, fungendo da giustificazione concettuale per le regole di inferenza.
